\subsubsection{偶遇性行为（一夜情 / 419）的安全守则：从"约"到"事后"}

上一节的术语词条回答了"这些词是什么意思"；但真正决定健康与安全结果的，是"具体怎么做"。本节把散落在全书各处的相关内容整合为一份可执行清单。\textbf{立场先行：成年人自愿的偶遇性行为本身不是道德问题，本文不做评判——这里只谈风险控制}，因为这类情境的风险结构与长期伴侣关系有明显不同：对方健康状况未知、关系无信任积累、事后往往无法追踪。

\vspace{0.5em}
\noindent\textbf{约定之前}：

\begin{itemize}
	\item \textbf{核验对方身份}：视频确认与照片一致再见面；把见面时间、地点、对方信息告知可信任的朋友（约会软件安全使用的完整指南见相关章节）
	\item \textbf{公共场所初见}：第一次见面选人多的公共场所；不轻易跟对方去私密空间，不去自己不掌握"退出权"的场所
	\item \textbf{饮品不离视线}：离开座位后不再饮用；警惕投放 GHB 等麻醉物导致的药物性侵（见网络性勒索与约会软件风险相关内容）
	\item \textbf{把期望提前说清}：是否发生性行为、是否全程使用安全套——见面前说清楚，比在床上僵持要体面得多；"只约会不性"（约素炮）同样是有效选项，说出口不丢人
\end{itemize}

\vspace{0.5em}
\noindent\textbf{过程中}：

\begin{itemize}
	\item \textbf{全程、全程、全程使用安全套}——包括口交（淋病、衣原体、梅毒、HPV 都可经口传播）。对方"上周刚检测过"不等于安全：检测存在窗口期（见"检测窗口期速查表"一节），窗口期内的阴性不能排除感染
	\item \textbf{同意随时可以撤回}：任何一方在任何时刻说停，就停——这同样适用于偶遇情境（同意的原则见"婚姻内的同意"一节，逻辑完全相同）
	\item \textbf{认清酒精与药物的角色}：严重醉酒状态下双方都无有效同意能力，且判断力与防护意愿都会下降；以助性药物"开party"叠加的化学性行为（chemsex）风险成倍放大（见物质使用相关章节）
	\item \textbf{不拍摄、不留影像}：任何"拍个照留念"都可能成为日后被要挟的把柄（见数字时代性隐私相关章节）
\end{itemize}

\vspace{0.5em}
\noindent\textbf{事后 72 小时与 2 周}：

\begin{itemize}
	\item \textbf{尽快排尿与清洁}：降低尿路感染风险（女性尤其重要）
	\item \textbf{若发生了无保护行为}：HIV 暴露后预防（PEP）\textbf{越早越好，72 小时内}仍可启动，最佳为 24 小时内——到当地传染病医院或疾控咨询；这是"事后可以补救的少数事情之一"，别拖
	\item \textbf{避孕失败的处理}：无保护阴道性行为后 72 小时内可服紧急避孕药（见"紧急避孕"一节）；注意它不能替代常规避孕
	\item \textbf{按窗口期安排检测}：4 周查 HIV/梅毒，12 周复查确认；淋病/衣原体等有症状随时查（见检测窗口期速查表）
	\item \textbf{接纳事后的情绪波动}：事后失落、空虚无感甚至莫名想哭都是可能的（性交后悲伤的相关机制见前文）——它与行为的对错无关，给情绪一点时间；也不必因此纠缠对方或自我指责
\end{itemize}

\vspace{0.5em}
\noindent\textbf{多伴侣与"炮友"场景的额外提醒}：参与者越多，每一次的感染风险越高，屏障规则需要更严格（每一次更换伴侣都更换安全套、手部清洁、玩具不混用）；固定"炮友"关系并不天然更安全——双方的"固定"不等于"互相唯一"，定期共同检测才是把风险管住的办法。

\begin{tcolorbox}[colback=pink!5!white,colframe=red!50!black,title=贴心小叮咛]
偶遇性行为的风险从来不是"道德风险"，而是"信息风险"——你不知道对方的健康状况，对方也可能不知道自己的。把本节的清单过一遍，尽可能把不可控变成可控；事后 72 小时内的那几通电话（PEP、紧急避孕、检测预约），永远比事后懊悔有用。
\end{tcolorbox}
